\section{\label{sec:future}Future Work}

The outstanding items of this work are several:

\begin{itemize} %TEMPORARY format

    \item Further refine the simulations, determining resolution and efficiency while employing more realistic conditions.

    \item Make a first-order conceptual design to gauge potential costs for a one-tonne demonstration prototype.

    \item Continue the now-ongoing investigation in the laboratory of a scintillator-filled quartz tube with \pmts\ at each end.
    This will be done at \uh\ in our detector laboratory, largely employing existing hardware.
    Our plans include outfitting parallel 1-m tubes with a source carriage controlled by a stepper motor.
    This will allow coincidence detection between the tubes.

    \item Investigate upscaling the detector to 100 tons.
    At present 10~000 tons seems out of reach, at least in the foreseeable future.

    \item Pursue the issues raised by our novel method of neutrino-direction extraction from the various configurations of IBD detectors described herein.

    \item Prepare a proposal for a modest demonstration detector.
    Without a sufficiently-powerful neutrino source available, we are investigating the utility of a neutron source and seeing double-scatter events between tubes.

    \item Study the number of events needed for discrimination of a 50-$\mathrm{MW_{th}}$ reactor in the field of much larger power reactors.

\end{itemize}
